\section{The distribution of \texorpdfstring{\(n\)}{n}-secants}
\label{sec:n-secants}

Let \(\mathcal K\) be a \((k,n)\)-arc in a projective plane of order \(q\).
We use the standard secant notation
\[
 \tau_i=\#\{\text{\(i\)-secants of \(\mathcal K\)}\},
\]
and write \(\rho_i(P)\) for the number of \(i\)-secants through
\(P\in\mathcal K\), and \(\sigma_i(Q)\) for the number through
\(Q\notin\mathcal K\).  Only \(i=n\) will occur below.  Put
\[
 b=q^2+q+1-k,\qquad \tau=q+1-n,\qquad T=\tau_n.
\]

The geometry supplies two bounds that need not hold for a selected block
family in an arbitrary design.

\begin{lemma}[Secants in a pencil]\label{lem:pencil-bounds}
\coverage{absent}
For every \(P\in\mathcal K\) and \(Q\notin\mathcal K\),
\begin{equation}\label{eq:pencil-bounds}
 \rho_n(P)\le D:=\left\lfloor\frac{k-1}{n-1}\right\rfloor,
 \qquad
 \sigma_n(Q)\le M:=\left\lfloor\frac{k}{n}\right\rfloor.
\end{equation}
Moreover,
\begin{align}
 \sum_{P\in\mathcal K}\rho_n(P)&=nT,&
 \sum_{Q\notin\mathcal K}\sigma_n(Q)&=\tau T,
 \label{eq:n-secant-incidences}\\
 \sum_{P\in\mathcal K}\binom{\rho_n(P)}2
 +\sum_{Q\notin\mathcal K}\binom{\sigma_n(Q)}2&=\binom T2.
 \label{eq:n-secant-pairs}
\end{align}
\end{lemma}

\begin{proof}
The \(n\)-secants through \(P\) have disjoint sets of \(n-1\) further
points of \(\mathcal K\).  Those through \(Q\) have disjoint \(n\)-point
intersections with \(\mathcal K\).  This proves \eqref{eq:pencil-bounds}.
The first two equations count incidences.  Every two distinct projective lines
meet once, either in \(\mathcal K\) or outside it, which gives
\eqref{eq:n-secant-pairs}.
\end{proof}

We will sometimes impose the stronger hypothesis
\begin{equation}\label{eq:lambda-coverage}
 \sigma_n(Q)\ge\lambda\qquad(Q\notin\mathcal K),
\end{equation}
where \(\lambda\ge1\).  For \(\lambda=1\), this is precisely completeness.
Rather than introduce another name for \eqref{eq:lambda-coverage}, we state
the multiplicity explicitly in each result.

\begin{corollary}[Integer feasibility for the \(n\)-secants]
\label{cor:n-secant-feasibility}
\coverage{absent}
Assume \eqref{eq:lambda-coverage}.  Then there is an integer \(T\) satisfying
\begin{gather}
 \left\lceil\frac{\lambda b}{\tau}\right\rceil
 \le T\le
 \left\lfloor\frac{k(k-1)}{n(n-1)}\right\rfloor,
 \label{eq:T-range}\\
 [\Phimin(b,\tau T),\Phimax(b,\lambda,M,\tau T)]
 \cap
 [\binom T2-\Phimax(k,0,D,nT),
  \binom T2-\Phimin(k,nT)]\ne\varnothing.
 \label{eq:arc-integer-condition}
\end{gather}
\end{corollary}

\begin{proof}
Equation~\eqref{eq:lambda-coverage} and
\eqref{eq:n-secant-incidences} give the lower bound for \(T\).  Each
\(n\)-secant contains \(\binom n2\) pairs of points of \(\mathcal K\), and
each pair lies on one line, giving the upper bound.  Apply
Theorem~\ref{thm:paired-envelope} to the point--line design of the plane,
using Lemma~\ref{lem:pencil-bounds} for the degree bounds.
\uses{thm:paired-envelope,lem:pencil-bounds}
\end{proof}

If
\[
 v_{\mathcal K}\equiv nT\pmod k,\qquad 0\le v_{\mathcal K}<k,
 \qquad
 v_{\mathrm{ext}}\equiv\tau T\pmod b,\qquad 0\le v_{\mathrm{ext}}<b,
\]
then the two minimum terms in \eqref{eq:arc-integer-condition} imply
\begin{equation}\label{eq:integer-incidence-bound}
 T\left(\frac{n^2}{k}+\frac{\tau^2}{b}-1\right)
 +\frac{v_{\mathcal K}(k-v_{\mathcal K})}{kT}
 +\frac{v_{\mathrm{ext}}(b-v_{\mathrm{ext}})}{bT}
 \le q.
\end{equation}
Discarding the two nonnegative residue terms gives the familiar real
incidence bound.  Since \(T\ge\lambda b/\tau\), we obtain
\begin{equation}\label{eq:classical-arc-bound}
 \frac{\lambda b}{\tau}
 \left(\frac{n^2}{k}+\frac{\tau^2}{b}-1\right)\le q.
\end{equation}
After denominators are cleared, \eqref{eq:classical-arc-bound} is the
quadratic in Bishnoi--Mattheus--Schillewaert~\cite[Theorem~8.1]{BMS}.
Thus the classical bound is the real relaxation of
\eqref{eq:arc-integer-condition}; the residue terms record part of the
discarded integer information.

\begin{remark}\label{rem:ordinary-completeness}
For a complete \((k,n)\)-arc one takes \(\lambda=1\).  The stronger values of
\(\lambda\) below mean only that every external point is on at least
\(\lambda\) \(n\)-secants.  This is the same multiplicity hypothesis used in
the dual form of Bishnoi--Mattheus--Schillewaert~\cite[Theorem~8.1]{BMS}.
\end{remark}
